% Осень26!A24; нумерация внутри столбца без пустых строк.
% Источник: лекция 3, §7 (строка 189)
\ticket{21}{Счетная аддитивность классического объема}
\label{ticket:21}

\begin{theorem}
Классический объём на полукольце ячеек в $\mathbb R^m$ является $\sigma$-конечной мерой.
\end{theorem}

\begin{proof}
Конечная аддитивность следует из разбиения ячеек общей сеткой: раскрытие произведения сумм длин рёбер даёт сумму объёмов частей.
По \hyperref[ticket:19]{критерию счётной аддитивности} достаточно проверить счётную полуаддитивность.

Пусть непустая ячейка $P=[a,b)$ покрыта ячейками $P_i=[a_i,b_i)$.
Для $\varepsilon>0$ выберем $a<b'<b$ так, что
\[
  P^-=[a,b'),\qquad \mu(P)-\mu(P^-)<\varepsilon.
\]
Для каждого $i$ выберем $a_i'<a_i$ так, чтобы
\[
  P_i^+=[a_i',b_i),\qquad
  \mu(P_i^+)<\mu(P_i)+\varepsilon2^{-i}.
\]
Открытые параллелепипеды $(a_i',b_i)$ покрывают компакт $[a,b']$.
Выберем конечное подпокрытие с множеством индексов $I$.
По конечной полуаддитивности
\[
  \begin{aligned}
  \mu(P)&<\mu(P^-)+\varepsilon
  \le\sum_{i\in I}\mu(P_i^+)+\varepsilon\\
  &\le\sum_{i=1}^{\infty}\mu(P_i)+2\varepsilon.
  \end{aligned}
\]
Ввиду произвольности $\varepsilon$ получаем требуемую полуаддитивность.
Наконец, целочисленные единичные ячейки счётно покрывают $\mathbb R^m$ и имеют конечный объём, что доказывает $\sigma$-конечность.
\end{proof}
